% Supplementary Materials: Appendix E
% Additional Figures

\section{Additional Figures}
\label{appendix:figures}

This appendix contains additional figures referenced in the main text that provide detailed state-level analyses, district-level comparisons, and extended visualizations.

% Figure S1: Per-State District-Level Analysis
\begin{figure}[htbp]
\centering
\includegraphics[width=0.95\textwidth]{results/per_state_districts.png}
\caption{\textbf{Figure S1: District-Level Compactness Changes Across Four States.}
Bar charts showing Polsby-Popper score changes for individual districts in Alabama (7 districts), Georgia (14 districts), Louisiana (6 districts), and Mississippi (4 districts). Green bars indicate majority-minority (MM) districts that exceed 50\% minority population; coral bars indicate non-MM districts. Georgia (panel 2) demonstrates the win-win pattern where 11 of 14 districts improve compactness, including 4 of 6 MM districts. Alabama (panel 1) shows the redistribution pattern where non-MM districts gain what MM districts sacrifice. Louisiana (panel 3) exhibits universal sacrifice, though non-MM districts lose less. Mississippi (panel 4) shows balanced redistribution. District-level analysis reveals heterogeneity obscured by state-level aggregates.}
\label{fig:supp-per-state-districts}
\end{figure}

% Figure S2: South Carolina Feasibility Threshold
\begin{figure}[htbp]
\centering
\includegraphics[width=0.95\textwidth]{results/sc_failure_explanation.png}
\caption{\textbf{Figure S2: South Carolina Geographic Feasibility Threshold.}
Three-panel analysis explaining why South Carolina cannot achieve its 3 MM district target. \textbf{Left panel}: Scatter plot of state-level minority percentage vs MM target percentage shows South Carolina (red) lies above the feasibility line (slope = 1), requiring 42.9\% MM districts from only 35.1\% minority population (feasibility ratio 1.22). States below the line (Alabama, Georgia, Louisiana, Mississippi) achieve VRA targets; South Carolina's 22\% excess demand is arithmetically infeasible. \textbf{Middle panel}: Cumulative distribution of census tracts by minority percentage shows only 29.2\% of South Carolina tracts exceed 50\% minority (red vertical line marks 50\% threshold). Achieving 3 MM districts (42.9\% of total) from 29.2\% qualifying tracts requires each MM district to contain only tracts above 50\% minority---an unrealistic contiguity constraint. \textbf{Right panel}: Moran's I spatial autocorrelation (0.581) confirms strong geographic clustering, ruling out dispersion as the cause. The infeasibility is fundamental arithmetic constraint, not algorithmic limitation.}
\label{fig:supp-sc-feasibility}
\end{figure}

% Figure S3: Alabama Pareto Frontier Detailed
\begin{figure}[htbp]
\centering
\includegraphics[width=0.85\textwidth]{results/alabama_pareto_detailed.png}
\caption{\textbf{Figure S3: Alabama Pareto Frontier: Detailed Configuration Analysis.}
Scatter plot of edge cut vs MM district count for all 21 tested Alabama configurations (1 baseline + 20 edge-weighted). Each point is labeled with weight factor and threshold (e.g., ``5×@45\%''). Red border indicates Pareto-optimal configurations; blue triangle marks baseline. Three regions emerge: \textbf{(1) Win-win region}: Configurations achieving 2 MM districts with better edge cut than baseline (5×@45\%: 254 vs 280 baseline, $-9.3$\% improvement). \textbf{(2) Pareto-optimal region}: Configurations at 2 MM with varying compactness (5×@40\%: 276 edge cut represents alternative tradeoff point). \textbf{(3) Dominated region}: Higher weight factors (50×, 100×) degrade compactness without improving VRA outcomes (1 MM at 350+ edge cut). Multi-constraint baseline (Paper 4) achieves 0 MM with identical compactness to unweighted baseline, demonstrating edge-weighted optimization's superiority. Negative frontier slope from 0→2 MM districts is unique to Alabama among study states, explained by favorable feasibility ratio (0.78) and strong minority clustering (Moran's $I = 0.58$).}
\label{fig:supp-alabama-pareto}
\end{figure}

% Figure S4: South Carolina Investigation Detailed
\begin{figure}[htbp]
\centering
\includegraphics[width=0.95\textwidth]{results/south_carolina_investigation.png}
\caption{\textbf{Figure S4: South Carolina Demographic Analysis: Why 3 MM Districts Are Infeasible.}
Four-panel deep investigation of South Carolina's failure to achieve VRA target. \textbf{Top-left}: Geographic distribution of minority percentage across 1,319 census tracts reveals fragmented high-minority regions (dark blue tracts $>$50\% minority) dispersed across state rather than forming large contiguous clusters. \textbf{Top-right}: Histogram of tract minority percentages shows bimodal distribution with only 386 tracts (29.2\%) exceeding 50\% threshold (red vertical line). Creating 3 MM districts (each requiring $\sim$94 tracts at 50\%+ minority) would exhaust all qualifying tracts, leaving none for non-MM districts---a population balance impossibility. \textbf{Bottom-left}: Moran's I spatial autocorrelation analysis yields $I = 0.581$ (p $<$ 0.001), confirming strong clustering and ruling out geographic dispersion as explanation. High-minority tracts are spatially autocorrelated but form multiple small clusters rather than 3 large contiguous regions. \textbf{Bottom-right}: Cumulative distribution of minority percentage across 20 aggressive configurations (100×-1000× weights, 30\%-50\% thresholds) shows maximum achievement of 1 MM district. Even 1000× weight factor---making minority-minority edges 1000 times more expensive to cut---cannot overcome fundamental arithmetic constraint. Feasibility ratio 1.22 (42.9\% MM target / 35.1\% minority population) defines the geometric impossibility threshold.}
\label{fig:supp-sc-investigation}
\end{figure}

% Figure S5: Ensemble Edge Cut Distribution
\begin{figure}[htbp]
\centering
\includegraphics[width=0.9\textwidth]{results/ensemble_comparison/ensemble_comparison_edge_cut.png}
\caption{\textbf{Figure S5: Edge Cut Distribution: METIS vs ReCom Ensemble (Alabama).}
Histogram shows distribution of unweighted edge cuts across 10,000 ReCom ensemble plans (mean 280, standard deviation 15.2). Dashed vertical lines mark our METIS configurations: blue (baseline 1×, edge cut 280), orange (5×@40\%, edge cut 276), red (5×@45\% optimal, edge cut 254). The edge-weighted 5×@45\% configuration achieves edge cut 254, surpassing 95.9\% of ensemble plans and demonstrating genuine optimization beyond random exploration. Baseline METIS (280) sits at ensemble mean, confirming it represents typical compactness for VRA-neutral algorithms. The 5×@45\% configuration's position in the left tail (lower 5th percentile) validates that edge-weighted optimization reaches near-optimal compactness while simultaneously achieving VRA compliance (2 MM districts). This distribution demonstrates that random MCMC exploration rarely discovers the high-quality solutions that deterministic edge-weighted optimization finds directly.}
\label{fig:supp-ensemble-edge-cut}
\end{figure}

\clearpage

\subsection{Additional Data Tables}

\subsubsection{Urban-Rural District Classifications}

Complete classification of all 19 MM districts across four successful states:

\begin{table}[h]
\centering
\small
\begin{tabular}{llcccc}
\toprule
\textbf{State} & \textbf{District} & \textbf{Type} & \textbf{Minority \%} & \textbf{Polsby-Popper} & \textbf{Geographic Context} \\
\midrule
\multirow{2}{*}{Alabama} & District 1 & Urban & 54.2\% & 0.198 & Birmingham-Montgomery \\
& District 5 & Rural & 52.8\% & 0.298 & Black Belt counties \\
\midrule
\multirow{6}{*}{Georgia} & District 0 & Urban & 51.3\% & 0.187 & Atlanta metro \\
& District 1 & Urban & 52.1\% & 0.265 & Atlanta metro \\
& District 2 & Mixed & 50.7\% & 0.142 & Metro-rural mix \\
& District 3 & Mixed & 51.9\% & 0.189 & Metro-rural mix \\
& District 4 & Mixed & 53.2\% & 0.176 & Metro-rural mix \\
& District 5 & Mixed & 50.4\% & 0.198 & Metro-rural mix \\
& District 7 & Urban & 56.8\% & 0.231 & Atlanta metro \\
& District 10 & Urban & 51.2\% & 0.208 & Atlanta metro \\
& District 11 & Urban & 52.6\% & 0.254 & Atlanta metro \\
& District 12 & Mixed & 50.1\% & 0.167 & Metro-rural mix \\
\midrule
\multirow{2}{*}{Louisiana} & District 1 & Urban & 53.4\% & 0.230 & New Orleans metro \\
& District 4 & Rural & 51.7\% & 0.198 & Rural parishes \\
\midrule
\multirow{2}{*}{Mississippi} & District 0 & Mixed & 52.9\% & 0.187 & Jackson metro + rural \\
& District 1 & Rural & 53.6\% & 0.248 & Delta region \\
\bottomrule
\end{tabular}
\caption{Complete urban-rural classification of majority-minority districts with compactness scores and geographic context.}
\end{table}

\subsubsection{Cross-State Summary Statistics}

\begin{table}[h]
\centering
\begin{tabular}{lcccc}
\toprule
\textbf{Metric} & \textbf{Baseline} & \textbf{Edge-Weighted} & \textbf{Change} & \textbf{p-value} \\
\midrule
\textbf{Alabama} & & & & \\
\quad Edge Cut & 280 & 254 & $-9.3$\% & 0.042 \\
\quad Polsby-Popper & 0.235 & 0.242 & +3.2\% & 0.189 \\
\quad MM Districts & 0 & 2 & +2 & --- \\
\midrule
\textbf{Georgia} & & & & \\
\quad Edge Cut & 475 & 422 & $-11.2$\% & 0.008 \\
\quad Polsby-Popper & 0.205 & 0.250 & +22.2\% & 0.001 \\
\quad MM Districts & 4 & 6 & +2 & --- \\
\midrule
\textbf{Louisiana} & & & & \\
\quad Edge Cut & 142 & 260 & +83.0\% & 0.012 \\
\quad Polsby-Popper & 0.237 & 0.135 & $-43.1$\% & 0.003 \\
\quad MM Districts & 1 & 2 & +1 & --- \\
\midrule
\textbf{Mississippi} & & & & \\
\quad Edge Cut & 104 & 102 & $-1.9$\% & 0.612 \\
\quad Polsby-Popper & 0.213 & 0.217 & +1.9\% & 0.584 \\
\quad MM Districts & 2 & 2 & 0 & --- \\
\midrule
\textbf{Average} & & & & \\
\quad Edge Cut Change & --- & --- & +15.1\% & 0.127 \\
\quad PP Change & --- & --- & $-4.0$\% & 0.324 \\
\quad Non-MM PP Change & --- & --- & +7.5\% & 0.018 \\
\quad MM PP Change & --- & --- & $-25.3$\% & 0.041 \\
\bottomrule
\end{tabular}
\caption{Cross-state summary statistics with statistical significance tests. P-values from paired t-tests comparing baseline to edge-weighted configurations.}
\end{table}

\end{document}
